\chapter{Preliminaries}
\label{chap:Preliminaries}

\section{Definitions}
\label{sec:Definitions}
This section fixes the notation and terminology used throughout the report. Throughout, $\mathcal{T}$ denotes the teacher and $T$ denotes the index of the final step of a trajectory.
 
\subsection{Agentic Trajectories}
 
A language model agent solves a task by interacting with an environment $\mathcal{E}$ that exposes a set of tools. Given a task prompt $x$, the agent alternates between generating text and receiving the results of the tool calls that text contains.
 
\begin{Definitions}[Agentic trajectory]
An agentic trajectory is a sequence
\[
  \tau = (s_0, a_0, o_0, s_1, a_1, o_1, \ldots, s_T, a_T),
\]
in which $s_t$ is the state at step $t$, $a_t$ is the action generated by the model, and $o_t$ is the observation returned by $\mathcal{E}$ in response to $a_t$. The initial state is the task prompt, $s_0 = x$, and states evolve by concatenation, $s_{t+1} = (s_t, a_t, o_t)$. The episode terminates at step $T$ when the model produces a final response or a turn limit is reached.
\end{Definitions}
 
The state $s_t$ is therefore the complete interaction history visible to the model at step $t$. The terms \emph{rollout} and \emph{episode} are used interchangeably with trajectory when the trajectory is sampled from the model being trained.
 
\begin{Definitions}[Agentic step]
An agentic step is a pair $(s_t, a_t)$, where $a_t$ is the maximal contiguous span of model-generated tokens between two consecutive environment observations. Every action other than the terminal action $a_T$ consists of optional free-text reasoning followed by exactly one tool call; the terminal action contains the final response.
\end{Definitions}
 
Under this Definitions, step boundaries are fixed by the structure of the interaction rather than by any property of the model's output distribution. Observation tokens $o_t$ are produced by the environment rather than by the policy, and are therefore excluded from the training loss. The adequacy of this boundary as a unit of supervision is examined in Chapter~3.
 
\subsection{Student, Teacher and Teacher Access}
 
\begin{Definitions}[Student policy]
The student is an autoregressive language model $\pi_\theta$ with trainable parameters $\theta$. For an action $a_t = (a_{t,1}, \ldots, a_{t,|a_t|})$ of $|a_t|$ tokens,
\[
  \pi_\theta(a_t \mid s_t) = \prod_{j=1}^{|a_t|} \pi_\theta\!\left(a_{t,j} \mid s_t, a_{t,<j}\right).
\]
\end{Definitions}
 
A tokenizer maps a string to a sequence of tokens drawn from a finite vocabulary $\mathcal{V}$. Models from different families generally use different tokenizers, so that the student vocabulary $\mathcal{V}_S$ and the teacher vocabulary $\mathcal{V}_{\mathcal{T}}$ differ and their next-token distributions are not defined over a common support.
 
\begin{Definitions}[White-box and black-box teacher access]
A teacher $\mathcal{T}$ is \emph{white-box accessible} if, for any context $c$, its next-token distribution $p_{\mathcal{T}}(\cdot \mid c)$ over $\mathcal{V}_{\mathcal{T}}$ can be read. It is \emph{black-box accessible} if it can be queried only as a sampling oracle: given a text prompt $c$, it returns generated text $y \sim \mathcal{T}(\cdot \mid c)$, and neither its logits, its vocabulary, nor its tokenizer are observable.
\end{Definitions}
 
Under black-box access, teacher and student interact only through text, and the cross-tokenizer problem does not arise. A single such interaction is termed a \emph{teacher query}, and the number of teacher queries per training step is used as the unit of cost in Chapter~4.
 
\subsection{On-Policy Distillation}
 
\begin{Definitions}[Off-policy and on-policy distillation]
In \emph{off-policy} distillation the student is trained on trajectories sampled from the teacher, $\tau \sim \mathcal{T}$. Supervised fine-tuning (SFT) on teacher demonstrations is the standard instance, minimising
\[
  \mathcal{L}_{\mathrm{SFT}}(\theta) = -\,\mathbb{E}_{\tau \sim \mathcal{T}} \Big[ \sum_{t=0}^{T} \log \pi_\theta(a_t \mid s_t) \Big].
\]
In \emph{on-policy} distillation the trajectories are sampled from the student itself, $\tau \sim \pi_\theta$, and the teacher supplies supervision on the states $s_t$ that the student actually visits.
\end{Definitions}
 
Off-policy training optimises the student only on the state distribution induced by the teacher. At inference time the student conditions on its own earlier actions, and a small deviation places it in states absent from its training data, where further errors become more likely. This mismatch is termed \emph{exposure bias}, and the resulting growth of error with trajectory length is termed \emph{compounding error}. On-policy distillation removes the mismatch by construction.
 
\begin{Definitions}[Divergences]
For distributions $p$ and $q$ over a common support, the Kullback--Leibler (KL) divergence is
\[
  D_{\mathrm{KL}}(p \,\|\, q) = \sum_{v} p(v) \log \frac{p(v)}{q(v)},
\]
and the Jensen--Shannon divergence is
\[
  D_{\mathrm{JS}}(p \,\|\, q) = \tfrac{1}{2} D_{\mathrm{KL}}(p \,\|\, m) + \tfrac{1}{2} D_{\mathrm{KL}}(q \,\|\, m), \qquad m = \tfrac{1}{2}(p + q).
\]
With the student as $p$ and the teacher as $q$, $D_{\mathrm{KL}}(\pi_\theta \,\|\, p_{\mathcal{T}})$ is termed the \emph{reverse} KL and $D_{\mathrm{KL}}(p_{\mathcal{T}} \,\|\, \pi_\theta)$ the \emph{forward} KL.
\end{Definitions}
 
Logit-based on-policy distillation minimises a per-token divergence of this kind between student and teacher along student-generated sequences. Because it requires $p_{\mathcal{T}}$ over a support shared with $\pi_\theta$, it requires both white-box access and a common vocabulary.
 
\subsection{Supervision Granularity and Credit Assignment}
 
\begin{Definitions}[Supervision granularity]
Supervision is \emph{trajectory-level} if it assigns a single scalar $r(\tau)$ to a complete trajectory, and \emph{step-level} if it assigns a sequence of scalars $\big(r(a_t \mid s_t)\big)_{t=0}^{T}$, one per agentic step. It is \emph{token-level} if it assigns a signal to every generated token, and \emph{chunk-level} if it assigns a signal to contiguous token spans chosen without reference to step boundaries.
\end{Definitions}
 
\begin{Definitions}[Credit assignment]
Credit assignment is the problem of attributing a trajectory's outcome to the individual actions $a_0, \ldots, a_T$ that produced it.
\end{Definitions}
 
Trajectory-level supervision leaves credit assignment entirely to the optimiser, which receives the same signal for every action in the trajectory regardless of which action was responsible for the outcome. Relatedly, an \emph{outcome reward} scores only the result of the episode, such as task success, whereas a \emph{process reward} scores intermediate steps. The step-level rewards studied in this thesis are process rewards derived from a black-box teacher.
 
\subsection{Rubric-Based Supervision}
 
\begin{Definitions}[Rubric]
A rubric is a finite set of weighted criteria $C = \{c_k\}_{k=1}^{K}$, where each item $c_k = (\rho_k, w_k)$ consists of a textual criterion $\rho_k$ and an importance weight $w_k > 0$. Each criterion must be decidable as satisfied or not satisfied for a single response in isolation.
\end{Definitions}
 
In Rubric-based On-Policy Distillation (ROPD) \cite{fang2026ropd}, rubrics are produced and applied by two components, both of which may be instantiated by the teacher itself.
 
\begin{Definitions}[Rubricator]
Given a prompt $x$, a set of $m$ teacher responses $Y^{\mathcal{T}}_x = \{y^{\mathcal{T}}_j\}_{j=1}^{m}$ and a set of $n$ student rollouts $Y^{S}_x = \{y^{S}_i\}_{i=1}^{n}$, a Rubricator $\mathcal{R}$ induces a rubric from the contrast between the two sets,
\[
  C_x = \mathcal{R}\big(x, Y^{\mathcal{T}}_x, Y^{S}_x\big).
\]
The rubric $C_x$ is shared across all student rollouts for the prompt $x$.
\end{Definitions}
 
\begin{Definitions}[Verifier and weighted pass rate]
A Verifier $\mathcal{V}$ returns a binary judgement $v_{i,k} = \mathcal{V}(x, y^S_i, c_k) \in \{0, 1\}$ indicating whether response $y^S_i$ satisfies criterion $\rho_k$. The Verifier is \emph{blind} if it is not told whether the response under evaluation was produced by the teacher or by the student. The reward of $y^S_i$ is its weighted pass rate
\[
  r(y^S_i) = \frac{\sum_{k=1}^{K} w_k \, v_{i,k}}{\sum_{k=1}^{K} w_k}.
\]
\end{Definitions}
 
The step-level analogue replaces the prompt $x$ with the state $s_t$ and whole responses with actions. Teacher reference actions are sampled at the student's own state, $A^{\mathcal{T}}_t = \{a^{\mathcal{T}}_{t,j}\}_{j=1}^{m}$ with $a^{\mathcal{T}}_{t,j} \sim \mathcal{T}(\cdot \mid s_t)$, a step-specific rubric $C_{s_t}$ is induced from the contrast between $A^{\mathcal{T}}_t$ and the student's action $a_t$, and the Verifier scores $a_t$ against $C_{s_t}$ to produce $r(a_t \mid s_t)$. The construction of this procedure is the subject of Chapter~3.
 
\begin{Definitions}[Rubric gaming]
Rubric gaming is the attainment of a high weighted pass rate by a response that does not exhibit the quality the rubric is intended to measure, for instance through formatting devices or keyword insertion. It is a special case of \emph{reward hacking}, in which a policy optimises a proxy reward at the expense of the objective the proxy was designed to represent.
\end{Definitions}
 
\subsection{Group Relative Policy Optimization}
 
Group Relative Policy Optimization (GRPO) \cite{shao2024deepseekmath} is the policy optimisation algorithm through which all rewards in this report are applied.
 
\begin{Definitions}[Group-relative advantage]
For a prompt $x$, GRPO samples a group of $n$ rollouts $\{y_i\}_{i=1}^{n}$ from the current policy $\pi_{\theta_{\mathrm{old}}}$ and obtains rewards $\{r_i\}_{i=1}^{n}$. The advantage of rollout $i$ is
\[
  A_i = \frac{r_i - \operatorname{mean}\big(\{r_j\}_{j=1}^{n}\big)}{\operatorname{std}\big(\{r_j\}_{j=1}^{n}\big) + \epsilon},
\]
where $\epsilon$ is a small constant for numerical stability.
\end{Definitions}
 
The advantage measures how much better a rollout performed than the other rollouts for the same prompt, which removes the need for a separately learned value function. With the per-token importance ratio
\[
  \rho_{i,j}(\theta) = \frac{\pi_\theta(y_{i,j} \mid x, y_{i,<j})}{\pi_{\theta_{\mathrm{old}}}(y_{i,j} \mid x, y_{i,<j})},
\]
GRPO maximises the clipped objective
\[
  J(\theta) = \mathbb{E}\!\left[ \frac{1}{n} \sum_{i=1}^{n} \frac{1}{|y_i|} \sum_{j=1}^{|y_i|} \min\!\Big( \rho_{i,j}(\theta) A_i,\ \operatorname{clip}\big(\rho_{i,j}(\theta), 1 - \epsilon_{\mathrm{clip}}, 1 + \epsilon_{\mathrm{clip}}\big) A_i \Big) \right] - \beta \, D_{\mathrm{KL}}\big(\pi_\theta \,\|\, \pi_{\mathrm{ref}}\big),
\]
where $\epsilon_{\mathrm{clip}}$ bounds the size of each policy update, $\pi_{\mathrm{ref}}$ is a fixed reference policy, and $\beta$ controls the strength of the KL penalty.
 
For agentic trajectories, the inner sum ranges over action tokens only, with observation tokens masked. In standard GRPO every token of rollout $i$ receives the same advantage $A_i$, which is precisely trajectory-level supervision. Step-level supervision replaces $A_i$ with a step-dependent advantage $A_{i,t}$ applied to the tokens of action $a_t$.
 

\section{Literature Review}
\label{sec:Literature Review}

This section surveys the work that motivates step-level, black-box on-policy
distillation for agentic tool-use tasks. Section~\ref{sec:lr-rlhf} covers the
post-training paradigm on which teacher-guided policy optimization rests.
Section~\ref{sec:lr-opd} introduces on-policy distillation and its two structural
assumptions. Sections~\ref{sec:lr-crosstok} and \ref{sec:lr-blackbox} cover the two
lines of work that relax those assumptions, and Section~\ref{sec:lr-multiteacher}
covers multi-teacher methods that address the single-teacher capability ceiling.
Section~\ref{sec:lr-gap} states the resulting gap.

\subsection{Post-training foundations}
\label{sec:lr-rlhf}

Modern language model post-training rests on the supervised fine-tuning, reward
modelling and reinforcement learning pipeline introduced for instruction following,
in which a reward model trained on human preference comparisons supplies the signal
for a subsequent PPO stage \autocite{ouyang_instructgpt_2022}. Direct Preference
Optimization later collapsed this pipeline into a single supervised-style loss over
preference pairs, dispensing with both the explicit reward model and the
reinforcement learning inner loop \autocite{rafailov_dpo_2023}. Both establish the
premise that later distillation methods inherit: a stronger signal, whether a learned
reward model or a teacher policy, steers the student policy rather than merely
supplying static labelled data.

Group Relative Policy Optimization (GRPO) replaces PPO's learned value function with
a group-normalised advantage computed across multiple sampled rollouts per prompt
\autocite{shao_deepseekmath_grpo_2024}. DeepSeek-R1 subsequently demonstrated that
large-scale reinforcement learning from a base model, without a supervised cold
start, is sufficient to induce extended chain-of-thought reasoning
\autocite{deepseekai_r1_2025}. GRPO forms the optimization backbone of most methods
discussed below. Black-box and multi-teacher distillation methods typically express
their teacher-derived signal as a reward inside a GRPO update rather than as a
separate objective.

Extending this paradigm to multi-turn, tool-using agents introduces credit-assignment
problems absent from single-turn reasoning
\autocite{agentic_rl_tooluse_2025a,agentic_rl_tooluse_2025b}. An agentic trajectory
interleaves model-generated text with external, non-differentiable tool outputs, and
a single early tool-call error can invalidate every subsequent step. This setting
motivates the remainder of the review, since it is where trajectory-level supervision
proves least adequate.

\subsection{On-policy distillation}
\label{sec:lr-opd}

On-policy distillation (OPD) departs from offline imitation of teacher demonstrations
by training the student on its own rollouts, with a teacher supplying dense
per-token correction on the states the student actually visits
\autocite{lu_thinkingmachines_opd_2025}. This addresses the exposure bias of
supervised fine-tuning, under which the student is optimized only on the teacher's
static state distribution and therefore encounters unfamiliar states at inference
time as small errors compound. Subsequent framing positions OPD as a distinct
post-training primitive alongside off-policy distillation and self-distillation, and
surveys its adoption across recent frontier model families
\autocite{xu_mopd_notion_2025}.

A recurring concern in this literature is that the gains available to OPD are bounded
by teacher quality, since a student cannot distill past what its teacher knows
\autocite{brown_opd_ceiling_2025}. This observation motivates the multi-teacher work
discussed in Section~\ref{sec:lr-multiteacher}. Two further assumptions in the
standard formulation are left largely unexamined. First, OPD as originally posed
requires read access to the teacher's token-level logits. Second, it applies
supervision as one continuous signal over a trajectory rather than as discrete
per-step credit. The following two sections cover work that relaxes each assumption
in turn.

\subsection{Cross-tokenizer distillation}
\label{sec:lr-crosstok}

A separate line of work relaxes the requirement that teacher and student share a
tokenizer, which is a precondition for direct logit matching whenever the two models
come from different families. The Universal Logit Distillation loss reformulates
matching as an optimal-transport distance between sorted teacher and student logit
distributions, removing the shared-vocabulary constraint in the offline setting
\autocite{boizard_uld_2025}. Dual-Space Knowledge Distillation projects teacher and
student representations into a shared space through cross-attention rather than
aligning logits directly \autocite{zhang_dskd_2024}, and MinED aligns mismatched
token sequences by greedy minimum-edit-distance matching \autocite{wan_mined_2024}.
Both remain offline and sequence-level.

Approximate Likelihood Matching (ALM) advances the field to a chunk-level, binarized
$f$-divergence formulation and forms the direct basis for GOLD. Its limitations
section names on-policy distillation as future work rather than as a solved problem
\autocite{minixhofer_alm_2025}. GOLD extends ALM and the Universal Logit
Distillation loss into the on-policy setting, correcting the sequence-truncation
behaviour of the latter and replacing indiscriminate logit sorting with a
direct-mapping-first vocabulary alignment strategy. It is implemented as
\texttt{GOLDTrainer} in the TRL library \autocite{patino_gold_2025}. GOLD nevertheless
remains a white-box method: it removes the shared-tokenizer requirement without
removing the requirement for teacher logits.

More recent work argues that these alignment strategies discard supervision precisely
at the token boundaries where teacher and student disagree \autocite{sun_simct_2026},
and proposes marginalizing over byte-level prefixes in place of committing to a
single alignment \autocite{byte_prefix_marginalization_2026}. Further background
includes multi-level optimal-transport formulations \autocite{multi_level_ot_2024},
a byte-level distillation interface \autocite{byte_level_interface_2026},
cross-tokenizer likelihood scoring, which separates the field into optimal-transport
and marginalized-likelihood camps \autocite{cross_tokenizer_likelihood_scoring_2025},
and CTKD \autocite{chen_ctkd_2025}.

The mechanism underlying this subfield is tokenization bias. Two statistically
equivalent tokenized and byte-level models may exhibit substantially different
predictive distributions, because canonical tokenization commits to a single
segmentation of a string and discards the alternative cover encodings over which a
byte-level model would implicitly marginalize \autocite{phan_tokenization_bias_2025}.
Tokenizer mismatch is therefore a source of genuine distributional divergence rather
than an implementation inconvenience.

Cross-tokenizer and black-box distillation address different constraints. The former
removes the shared-vocabulary requirement while retaining logit access; the latter
removes logit access outright. Both respond to the same practical circumstance, in
which the strongest available teachers are proprietary models that expose neither
tokenizer nor logits.

\subsection{Black-box on-policy distillation}
\label{sec:lr-blackbox}

Black-box OPD removes the logit-access requirement, deriving supervision from the
teacher's generated text alone. Generative Adversarial Distillation (GAD) trains an
auxiliary discriminator under a Bradley-Terry pairwise loss, co-evolving with the
student, to score complete student responses; the resulting scalar is optimized as a
GRPO reward \autocite{ye_gad_2025}. The signal is one scalar per finished response,
and is therefore no denser than the outcome-based reinforcement learning that OPD was
introduced to improve upon. The reported results are consistent with a discriminator
that tracks surface style rather than correctness: AIME24 accuracy moves from 24.17\%
to 27.52\%, and evaluation is confined to open-ended chat, with no reasoning or
agentic benchmarks.

Rubric-based On-Policy Distillation (ROPD) induces a prompt-specific weighted rubric
from a contrast between teacher and student responses. A verifier scores the student
rollout against the rubric blind to response provenance, and the weighted pass rate
becomes the GRPO reward \autocite{fang_ropd_2026}. ROPD reports cross-architecture
gains on mathematics, science and medicine benchmarks without tokenizer alignment,
together with up to a tenfold improvement in sample efficiency over logit-based
baselines, and observes rubric gaming in fewer than 2\% of early-training rollouts.
Its supervision nevertheless remains trajectory-level: a single scalar computed after
a complete response and applied retrospectively through GRPO. Rubric induction also
requires two teacher calls per rollout at every training step.

OmniOPD removes logit access at a granularity finer than the trajectory. It selects
high-predictive-entropy chunks of 50 to 100 tokens from the student's chain of
thought, samples $N$ Monte Carlo continuations from the black-box teacher conditioned
on each chunk prefix, and scores the student's chunk against them under a continuous
semantic similarity metric such as normalised edit distance or ROUGE-1 unigram
overlap \autocite{zhou_omniopd_2026}. A Dirichlet-Multinomial prior bounds estimator
variance and a KL trust region constrains policy drift on unaudited tokens. Across
mathematical reasoning and competitive programming benchmarks, with Qwen3,
Claude-4.5-Haiku and Gemini-2.5-Flash as teachers, OmniOPD exceeds both supervised
fine-tuning and white-box OPD in most configurations. Its unit of supervision is a
contiguous token span within a single continuous chain of thought, anchored to the
student's own predictive entropy rather than to any structural boundary in the
trajectory. No tool call, environment turn, or agentic step is represented in either
its chunking mechanism or its evaluation suite.

\subsection{Multi-teacher distillation and the capability ceiling}
\label{sec:lr-multiteacher}

Orthogonal to teacher access is the observation that any single-teacher method is
capped by that teacher's own capability, since an erring teacher propagates its error
to the student. MAD-OPD recasts the teacher as a deliberative collective: $K$ teachers
debate over two rounds on the student's on-policy state, and each teacher's
contribution to the resulting token-level supervision is weighted by its post-debate
confidence \autocite{wang_madopd_2026}. A 4B student trained under a 14B and 8B
teacher debate surpasses its own 14B teacher on LiveCodeBench, which constitutes the
empirical demonstration that debate-derived supervision can exceed the single-teacher
ceiling.

MAD-OPD is also the only work surveyed here that extends OPD to agentic tasks at the
step level. Its On-Policy Agentic Distillation (OPAD) component introduces step-level
sampling by force-decoding, in order to stabilise training against the error
compounding that per-step mistakes induce over long agentic trajectories. The paper
further grounds its choice of divergence theoretically, proving that the per-token
gradient of reverse KL is unbounded under the privileged gap that debate induces,
whereas that of Jensen-Shannon divergence is bounded by 2. This motivates
Jensen-Shannon divergence for agentic tasks and reverse KL for code, where the
privileged gap arises to a lesser degree.

The method is nevertheless white-box and shared-vocabulary only. OPAD's step-level
sampling operates on token-level logits from the debating teachers, requiring both
logit access and a common vocabulary across the teacher ensemble and the student. The
authors name black-box and cross-vocabulary distillation as open extensions rather
than as problems the paper addresses \autocite{wang_madopd_2026}.

\subsection{The gap}
\label{sec:lr-gap}

The preceding sections trace two independent relaxations of standard OPD along two
axes: teacher access, from white-box to black-box, and supervision granularity, from
trajectory-level to step-level. Table~\ref{tab:lr-gap} positions existing methods
along these axes for agentic tool-use trajectories.

\begin{table}[H]
\centering
\begin{tabular}{|l|c|c|}
\hline
 & \textbf{Trajectory-level} & \textbf{Step-level} \\
\hline
\textbf{White-box (logits)} & GKD, logit-based OPD & OPAD (MAD-OPD) \\
\hline
\textbf{Black-box (text only)} & ROPD & Open \\
\hline
\end{tabular}
\caption{On-policy distillation methods positioned by teacher access and supervision
granularity for agentic tool-use trajectories.}
\label{tab:lr-gap}
\end{table}

No existing method combines black-box teacher access with step-level credit
assignment over multi-turn, tool-using agentic trajectories. ROPD achieves black-box
access at trajectory granularity alone. OPAD achieves step-level granularity for
agentic tasks but requires white-box, shared-vocabulary access. OmniOPD achieves
black-box access below trajectory granularity, but defines its chunks over continuous
chain-of-thought tokens rather than agentic steps, and is evaluated on single-turn
reasoning tasks only.

The open cell is the subject of this thesis. The question is whether rubric-based,
blind-verifier supervision of the kind ROPD applies to whole responses can instead be
applied per agentic step, and whether doing so improves over trajectory-level
black-box supervision on multi-turn tool-use tasks. Closing this cell is additionally
a precondition for transferring the multi-teacher debate mechanism of MAD-OPD into
black-box, cross-vocabulary space, where debate-derived step-level supervision could
be drawn from teachers spanning different model families without the shared-vocabulary
constraint OPAD currently imposes. That extension lies beyond the present scope.

\section{Contributions}
\label{sec:Contributions}


